\section{La posición de SGLang: la capa intermedia en la arquitectura de cinco capas de la IA}

Zhao Chenyang (赵晨阳) comienza presentándose: licenciatura en ciencias de la computación en Tsinghua, egresado en 2024, cursó quince meses de doctorado en UCLA y después, alrededor de noviembre de 2025, comenzó a emprender, uniéndose a una empresa orientada a motores de inferencia. El presentador dirige enseguida la pregunta hacia SGLang: cómo pasó de ser una organización de código abierto a convertirse en empresa, y qué experiencia tuvo Zhao Chenyang en ese proceso.

Primero explica qué es un motor de inferencia. Cuando un usuario común entra en contacto con la IA, lo que ve son productos de frontend como ChatGPT, Claude, AI PPT, Sora o AI Coding; pero una vez que la solicitud del usuario llega al backend, cómo el servidor programa el modelo, cómo administra el KV cache, cómo cuantiza, cómo maneja distintas formas de modelo, cómo aumenta el rendimiento y reduce la latencia: eso es lo que debe resolver el motor de inferencia. El \texttt{model.generate} de Hugging Face ayuda a entender el proceso de generación más elemental, pero atender solicitudes reales en producción exige una gran cantidad de optimización a nivel de sistema.

Zhao Chenyang recurre a una arquitectura de cinco capas para ubicar la inferencia: en la capa superior están los agentes y las aplicaciones; debajo, los modelos concretos, como Claude, GPT o modelos de código abierto; más abajo, los motores de inferencia, como OpenAI internal engine, TensorRT-LLM o SGLang; más abajo, el hardware, y más abajo aún, infraestructura como los semiconductores y la energía eléctrica. Cuanto más arriba, más rápido cambia; cuanto más abajo, más lento es el ciclo; el motor de inferencia se ubica justo en medio, como la capa de enlace clave entre los fabricantes de modelos y los fabricantes de hardware.

\begin{knowledgebox}{Por qué importa el motor de inferencia}
\begin{tabularx}{\textwidth}{>{\raggedright\arraybackslash}p{0.24\textwidth} Y}
\toprule
Nivel & Problema que debe resolver el motor de inferencia \\
\midrule
Programación de solicitudes & Cómo poner en cola de manera eficiente a distintos usuarios, distintas longitudes de contexto y distintos lotes. \\
KV Cache & En contextos largos y conversaciones de varios turnos, cómo reutilizar, trasladar y liberar la caché. \\
Forma del modelo & Procesos de inferencia distintos, como autorregresivo, difusión o flow matching, tienen necesidades de sistema distintas. \\
Estructura de costos & La latencia, el rendimiento, la memoria de video y el aprovechamiento del hardware determinan en conjunto si el producto puede escalar. \\
Adaptación del ecosistema & Los modelos de la capa superior cambian rápido y el hardware de la capa inferior cambia lento; la capa intermedia debe absorber continuamente los cambios de ambos lados. \\
\bottomrule
\end{tabularx}
\end{knowledgebox}

Lo más valioso que hay que conservar de este pasaje no son los detalles técnicos concretos de SGLang, sino el juicio de Zhao Chenyang sobre la «capa intermedia». En la era de la IA, cuanto más arriba está la capa, más ruidosa es y más fácil que la llene la narrativa del producto; pero lo que en realidad determina si una aplicación puede escalar suele ser el costo, el rendimiento y la confiabilidad de ingeniería de la capa intermedia. El motor de inferencia no genera noticias con la facilidad con que lo hace el lanzamiento de un modelo, ni enfrenta directamente al usuario como una aplicación, pero se convierte en un cimiento importante para que funcione toda la economía de la IA.

\subsection{Resumen de la sección}

La capa de motores de inferencia en la que se ubica SGLang es la capa intermedia entre las aplicaciones de IA, los modelos y el hardware. No responde directamente «qué tan inteligente es el modelo», sino «cómo mantener el servicio del modelo funcionando de manera estable con un costo asumible». La experiencia de Zhao Chenyang es adecuada para discutir la ola de abandono del doctorado precisamente porque no participó en un proyecto académico abstracto, sino en un sistema de código abierto situado en el centro de la estructura de costos de la industria.
